% © 2026 Ing. David Jaroš — CC BY-NC-ND 4.0
% UBT-AI-PROVENANCE-BEGIN
% schema: ubt-ai-provenance/v1
% tier: B_machine_verified
% ai_assistance: disclosed
% human_review: machine-verification
% editorial_responsibility: Ing. David Jaroš
% policy: ../../AI_PROVENANCE.md
% notice: Machine-verified against named sources or verifiers; individual attestation is not claimed.
% UBT-AI-PROVENANCE-END

\documentclass[12pt,a4paper]{article}
\usepackage[a4paper,margin=2.5cm]{geometry}
\usepackage{amsmath,amssymb,amsthm,mathtools}
\usepackage[most]{tcolorbox}
\usepackage[colorlinks=true,linkcolor=blue]{hyperref}
\usepackage{ubtprovenance}
\UBTTier{B}
\newtheorem{definition}{Definition}[section]
\newtheorem{theorem}[definition]{Theorem}
\newtheorem{corollary}[definition]{Corollary}
\theoremstyle{remark}\newtheorem{remark}[definition]{Remark}
\title{GAP-10T: Algebraic Torsion Selection in the Minimal First-Order Branch}
\author{Ing. David Jaro\v{s}}
\date{16 July 2026}
\begin{document}
\maketitle
\UBTProvenanceNotice

\begin{abstract}
The kinematic identity $\omega=\mathring\omega(e)+K(T)$ does not determine
$T$.  This note closes a sharply defined dynamical subgap.  In the minimal
Hilbert--Palatini/Einstein--Cartan branch, with nondegenerate tetrad and no
independent torsion kinetic term, variation of the Lorentz connection gives an
algebraic Cartan equation.  Its torsion map is pointwise invertible: vanishing
spin current forces $T=0$, while a specified spin current determines a unique
contorsion.  For a two-sided UBT matter derivative the spin current is obtained
by varying the same paired connection acting on $\Theta$.  What remains open is
not the algebra of torsion elimination, but the derivation of this minimal
first-order action, its exact UBT representation and normalization, and the
resulting $\Theta$ spin current from the fundamental theory.
\end{abstract}

\section{First-order branch}
Let $e^a$ be a nondegenerate tetrad, $\omega^{ab}=-\omega^{ba}$ an independent
Lorentz connection, and
\[
 R^{ab}=d\omega^{ab}+\omega^a{}_c\wedge\omega^{cb},\qquad
 T^a=de^a+\omega^a{}_b\wedge e^b.
\]
Consider the minimal first-order action
\begin{equation}
 S_{\rm HP}[e,\omega]
 =\frac{1}{4\kappa}\int
 \epsilon_{abcd}\,e^a\wedge e^b\wedge R^{cd}(\omega)
 -\frac{\Lambda}{24\kappa}\int
 \epsilon_{abcd}\,e^a\wedge e^b\wedge e^c\wedge e^d
 \label{eq:HP}
\end{equation}
plus a matter/UBT action $S_{\rm m}[e,\omega,\Theta]$.  Define the Lorentz-spin
current three-form by the convention
\begin{equation}
 \delta_\omega S_{\rm m}
 =\frac12\int \tau_{ab}\wedge\delta\omega^{ab}.
 \label{eq:spin-current-def}
\end{equation}

\section{Connection variation}
\begin{theorem}[Cartan connection equation]
For compactly supported variations of $\omega$, the connection Euler--Lagrange
equation of Eq.~\eqref{eq:HP} is
\begin{equation}
 \boxed{\epsilon_{abcd}\,T^a\wedge e^b=\kappa\,\tau_{cd}.}
 \label{eq:cartan}
\end{equation}
The equation is algebraic in torsion: it contains no derivative of $T$.
\end{theorem}

\begin{proof}
Use $\delta R^{cd}=D_\omega\delta\omega^{cd}$ and integrate by parts.  Since
$D_\omega(e^a\wedge e^b)=T^a\wedge e^b-e^a\wedge T^b$, antisymmetry of
$\epsilon_{abcd}$ combines the two terms into
$2\epsilon_{abcd}T^a\wedge e^b$.  Adding
Eq.~\eqref{eq:spin-current-def} gives Eq.~\eqref{eq:cartan}.  The cosmological
term is independent of $\omega$.
\end{proof}

\begin{theorem}[Pointwise invertibility of the Cartan torsion map]
For a nondegenerate tetrad in four dimensions, the linear map
\begin{equation}
 \mathcal C_e:\quad T^a\longmapsto
 \epsilon_{abcd}\,T^a\wedge e^b
 \label{eq:cartan-map}
\end{equation}
from the 24 torsion components to the 24 Lorentz-valued three-form components
is an isomorphism.
\end{theorem}

\begin{proof}
The claim is pointwise and Lorentz covariant, so choose an orthonormal frame
$e^a=dx^a$.  Expanding the six antisymmetric pairs $(c,d)$ in the four
three-form basis elements gives a $24\times24$ linear system.  Direct row
reduction has rank 24; the exact verifier is
\texttt{tools/verify\_remaining\_gr\_subclosures.py}.  Hence the kernel is
zero in one orthonormal frame and therefore in every nondegenerate frame.
\end{proof}

\begin{corollary}[Vacuum/minimal bosonic torsion]
If the matter action has zero Lorentz-spin current,
$\tau_{ab}=0$, then Eq.~\eqref{eq:cartan} uniquely gives
\begin{equation}
 \boxed{T^a=0,\qquad K^{ab}=0,\qquad
 \omega^{ab}=\mathring\omega^{ab}(e).}
\end{equation}
\end{corollary}

\begin{corollary}[Spin-sourced torsion]
For a specified $\tau_{ab}$, Eq.~\eqref{eq:cartan} determines a unique torsion
and hence a unique contorsion $K(T)$.  In the minimal branch torsion is
nonpropagating: it is an auxiliary field algebraically tied to spin current.
\end{corollary}

\section{Two-sided UBT spin current}
For the generic algebra-native derivative
\begin{equation}
 D_\mu\Theta=\partial_\mu\Theta+A_\mu\Theta-\Theta B_\mu,
\end{equation}
a Lorentz variation acts through a paired generator
\begin{equation}
 \mathcal J_{ab}\Theta
 =J^{L}_{ab}\Theta-\Theta J^{R}_{ab}.
\end{equation}
If the kinetic term uses an invariant real bilinear
$\langle\cdot,\cdot\rangle$, variation of the connection produces a current
schematically of the form
\begin{equation}
 \tau_{ab}\ \propto\
 *\,\operatorname{Sym}\langle D\Theta,\mathcal J_{ab}\Theta\rangle,
 \label{eq:theta-spin-current}
\end{equation}
with the exact sign and normalization fixed by the canonical action and its
left/right representation.  Equation~\eqref{eq:theta-spin-current} is not set
to zero by calling $\Theta$ ``bosonic'': a field carrying a nontrivial Lorentz
representation can source the independent spin connection.

\section{What is closed and what remains}
\begin{tcolorbox}[colback=green!6!white,colframe=green!55!black,title=Defensible status]
\textbf{GAP-10T-PALATINI: CLOSED CONDITIONALLY [L1].}
For the minimal Hilbert--Palatini branch with nondegenerate tetrad and no
independent torsion kinetic term, the connection equation is algebraic and
pointwise invertible.  Zero spin current gives zero torsion; specified spin
current gives unique contorsion.

\medskip
\textbf{GAP-10T-DYN: NARROWED, NOT FULLY CLOSED.}
The remaining UBT task is to derive the first-order gravitational term and the
precise paired matter representation from the canonical action, calculate its
spin current and normalization, and decide whether additional torsion-squared
or propagating-torsion invariants are present.
\end{tcolorbox}

The standard first-order geometry used here is not claimed as a new result.
The UBT-specific gain is the exact separation between (i) the already closed
kinematic reconstruction $\omega(e,T)$, (ii) the algebraically closed minimal
torsion equation, and (iii) the still-open origin and field content of the
canonical UBT action.

\begin{thebibliography}{9}
\bibitem{hehl1976} F.~W.~Hehl, P.~von der Heyde, G.~D.~Kerlick, and J.~M.~Nester,
``General Relativity with Spin and Torsion: Foundations and Prospects,''
\emph{Rev. Mod. Phys.} \textbf{48}, 393--416 (1976).
\bibitem{wald1984} R.~M.~Wald, \emph{General Relativity}, University of Chicago Press (1984).
\end{thebibliography}
\end{document}
